\subsection{Multiplikation mit einem Skalar}
\label{subsec:vector_algebra_scalar_multiplication}

Eine einzelne Zahl heißt in der Vektorrechnung \emph{Skalar}
\EnglishTerm{scalar}. Sie wirkt gleichzeitig auf alle Komponenten eines
Vektors.

\begin{DefinitionBox}{Skalarmultiplikation}{vector_algebra_scalar_multiplication}
  Für $\lambda\in\R$ und $\vec a\in\R^n$ gilt
  \[
    \lambda\vec a
    =\lambda\begin{pmatrix}a_1\\a_2\\\vdots\\a_n\end{pmatrix}
    =\begin{pmatrix}\lambda a_1\\\lambda a_2\\\vdots\\\lambda a_n\end{pmatrix}.
  \]
\end{DefinitionBox}

\begin{center}
  \CachedTikzFigure{content/resources/chapters/vector_algebra/tikz/scalar_multiplication_cases.tex}
\end{center}

\begin{KeyPointBox}{Geometrische Wirkung eines Skalars}{vector_algebra_scalar_effect}
  \[
    \begin{array}{ccl}
      |\lambda|>1
      &\Longleftrightarrow&\text{Streckung},\\[0.45em]
      0<|\lambda|<1
      &\Longleftrightarrow&\text{Stauchung},\\[0.45em]
      \lambda<0
      &\Longleftrightarrow&\text{zusätzliche Richtungsumkehr},\\[0.45em]
      \lambda=0
      &\Longleftrightarrow&\lambda\vec a=\vec 0.
    \end{array}
  \]
\end{KeyPointBox}

\begin{TheoremBox}{Rechenregeln der Skalarmultiplikation}{vector_algebra_scalar_laws}
  Für $\alpha,\beta\in\R$ und $\vec a,\vec b\in\R^n$ gilt
  \[
    \begin{aligned}
      \alpha(\beta\vec a)&=(\alpha\beta)\vec a,\\[0.4em]
      (\alpha+\beta)\vec a&=\alpha\vec a+\beta\vec a,\\[0.4em]
      \alpha(\vec a+\vec b)&=\alpha\vec a+\alpha\vec b.
    \end{aligned}
  \]
\end{TheoremBox}

\begin{ExampleBox}{Einen Vektor skalieren}{vector_algebra_scale_example}
  Für $\vec a=(2,-1,3)^{\mathsf T}$ gilt
  \[
    -2\vec a
    =-2\begin{pmatrix}2\\-1\\3\end{pmatrix}
    =\Cloze{3cm}{\begin{pmatrix}-4\\2\\-6\end{pmatrix}}.
  \]
  Der Faktor $-2$ verdoppelt die Länge und kehrt die Richtung um.
\end{ExampleBox}
